\section{Features}
\label{sec:features}

A tour of what the app does, grouped by what someone is trying to get done rather than by code
package. The navigation graph declares fourteen destinations: five bottom-navigation tabs (Home,
Browse, Search, AI, Settings) and nine pushed screens with typed arguments, each of which gets its
turn below.

\subsection{The knowledge base}

Branches, folders and posts are three values of one \code{type} column in one self-referencing
\code{nodes} table, so create, rename, move, reorder and delete are written once and behave the
same at every level. \S\ref{sec:data} covers the schema.

\textbf{Browse} is the tree walker. A breadcrumb row starts at \emph{Library} and adds a tappable
crumb per level, so jumping back three folders is one tap; system back goes up one level. The
\code{+} button offers \emph{Branch}, \emph{Folder} or \emph{Post}, defaulting to Branch at the root
and Post elsewhere, and choosing Post opens the editor with the current folder as the parent.

Each row has an overflow menu with \emph{Rename}, \emph{Move} and \emph{Delete}. Dragging reorders
among siblings, written when the drag ends; swiping left confirms, then deletes. \emph{Move} is a
computed list dialog rather than a drag --- the deliberate limit is that dragging reorders but never
re-parents --- and \code{possibleParents} offers \emph{Library (top level)} plus every branch and
folder that is not inside the item being moved.

The rules the tree enforces are few, and most exist to keep something else honest:

\noindent
\begin{tabularx}{\linewidth}{L{4.6cm} X}
\toprule
\textbf{Rule} & \textbf{Why} \\
\midrule
Only branches and folders take children &
A child of a post would be invisible in Browse \\
\addlinespace
A move that would cycle is refused &
\code{TreeRules.wouldCreateCycle} walks the parent chain before the write, so a folder cannot land
inside its own subtree \\
\addlinespace
Two posts cannot share a title &
Titles address posts everywhere --- \code{[[wiki-links]]}, AI actions, lookup by name --- so a
duplicate makes every lookup arbitrary. Checked only when a title changes, so an old import with
duplicates stays editable \\
\addlinespace
Deleting a container deletes its contents &
The foreign key cascades, the confirmation says so in words, and there is no undo \\
\bottomrule
\end{tabularx}

\medskip
A first launch is not an empty screen: Room's database-created callback seeds a \emph{Getting
Started} branch --- a welcome post, a \emph{How this library works} folder with two more, three
tags, one per-post glossary term and one global one, wiki-linked so backlinks, related posts and
the graph have something to show on day one. Seeding is tied to the database file being created,
not to the library being empty: the seed text invites deletion, and re-growing it behind the
reader's back would be a bug.

\subsection{Reading and writing}

The \textbf{reader} renders Markdown through Markwon, entirely on-device. Above the article sit
the title, the breadcrumb path, a star, a status chip and buttons for \emph{Edit} and \emph{Show in
graph}; below it come glossary chips, references, \emph{Linked from} and \emph{Related}, each
hidden when empty, so a short post is a short screen. Opening a post bumps its \code{updatedAt}, so
\emph{Recent posts} reflects reading and not only editing.

\code{[[Wiki-links]]} are resolved at render time rather than stored. Before Markwon sees the text,
each \code{[[Title]]} becomes a Markdown link on an internal scheme: \code{mocha://post/<id>} when
the title matches a post, \code{mocha://missing/<title>} when it does not. A custom resolver
intercepts both --- the first navigates, the second reports ``\emph{Title} is not in your library
yet''. An unresolved link is visible and inert rather than an error, because linking to a post you
have not written is a normal way to plan.

The \textbf{editor} is a title field, a \emph{Source}/\emph{Preview} tab pair and a toolbar: bold,
italic, heading, list, external link, wiki link, glossary term, add reference. Formatting buttons
wrap the selection or insert a prefix at the start of the line. The wiki-link button wraps a
selection in brackets, or with nothing selected lists every other post title and inserts the one
picked. \emph{Preview} runs the reader's own rewrite, so links look exactly as they will when saved.

Every keystroke goes to the ViewModel and is rendered back from state, so rotation cannot lose a
draft, and leaving with unsaved changes asks first. Tags are typed as a comma-separated list and
deduplicated case-insensitively. References and glossary terms are staged in the draft: one
transaction on save updates the node, rebuilds its outgoing link rows from the markdown, replaces
its tag rows and upserts its terms.

Renaming a post is the interesting write. \code{KnowledgeSync.retitle} writes the new title, then
rewrites \code{[[old title]]} in every post that links to it, in the same transaction, matching
case-insensitively and tolerating padding inside the brackets. One routine backs the rename in
Browse and in the editor. Without it an inbound link would degrade to ``not created yet'' and lose
its link row on the linking post's next save.

Deliberate limits: Markdown only --- no images, no attachments, no rich-text toolbar --- and the
source pane is a plain text field with no syntax highlighting.

\subsection{Connecting knowledge}

Three deliberately different answers to ``what else is this connected to'' sit under each article.
\textbf{Linked from} is exact: the posts whose markdown contains a resolved wiki-link to this one,
read from the \code{links} table that saving maintains.

\textbf{Related} is inferred. It scores every other post and shows the top five: $+2$ per tag
shared with this post; $+2$ for each post this one links out to; and $+1$ for each post that also
links to one of those targets --- co-citation, so two posts that both cite \emph{Raft} surface for
each other with no direct link between them. Posts already listed under \emph{Linked from}, and the
post itself, are excluded before scoring.

The \textbf{knowledge graph} draws the same data as a picture. Posts are dots sized by degree and
filled by learning status, favourites accented; edges are wiki-links, and a chip switches between
\emph{Links only} and \emph{Links + tags}, the second adding an edge between posts sharing a tag.
From Home the graph is the whole library; from a post's \emph{Show in graph} it is that post plus
its one-hop neighbourhood, which a \emph{Whole library} chip widens. Two limits are made visible
rather than hidden: past 250 posts only the most-connected are drawn and the caption says so
(``\emph{n} most-connected of \emph{m} posts''), and a tag on more than eight posts is drawn as a
ring rather than a clique, because $O(k^2)$ grey edges would bury the wiki-links.

The layout is the project's own code: \code{util/ForceLayout.java}, about 150 lines of
Fruchterman--Reingold in framework-free Java --- repulsion between every pair, attraction along
edges, a cooling temperature, sunflower-spiral seeding so isolated nodes start spread out instead
of piled on one radius, and a fixed seed so the same library draws the same shape between launches.
It returns a flat \code{float[]} in a fixed square world that \code{GraphView} fits to whatever
screen it gets, so rotation never re-runs the layout. Layout runs on \code{Dispatchers.Default};
the view handles drag-to-pan, pinch-to-zoom, double-tap-to-fit, tap-to-select and tap-again-to-open,
and hides node labels when the dots are packed too tightly to read.

\subsection{Learning metadata}

A post carries four pieces of state beyond its text, each earning its place by feeding another
screen. \textbf{Status} is \emph{None}, \emph{Reading}, \emph{In progress} or \emph{Finished}, set
from a chip in the reader or the editor. It is a reading state, not a memory schedule --- there is no
spaced repetition --- and it drives \emph{Continue reading} on Home (posts in \emph{Reading} or
\emph{In progress}, most recently touched first) and two search filters. \textbf{Favourite} is the
star in the reader header, filling the Home section, the Favorites screen, the favourites search
filter and the accent ring in the graph. \textbf{Tags} are stored as shared rows; they appear as
chips under the reader title, tapping one opens a screen of the posts carrying it, and a
\textbf{Tags} screen lists every tag alphabetically with a post count and a filter box.

\textbf{References} are external material: a kind (YouTube, Article, Book, Other), an optional
title and a URL, attached from the editor and shown as cards under the article captioned with the
link's host. A post's reference list is the union of the rows someone added on purpose and the
YouTube URLs sitting in the markdown right now, detected at read time and deduplicated by video id.
The derived half is not stored, which stops a re-save from deleting a reference the user or the
assistant added by hand; the trade is that a derived card is removed by editing the markdown rather
than by a close button, and the editor marks it accordingly.

Tapping a YouTube card plays it in a bottom sheet without leaving the article: a locked-down
\code{WebView} loading the IFrame embed against \code{youtube-nocookie.com}, file and content
access off, navigation restricted to an allow-list of YouTube hosts and everything else handed to
the browser. An \emph{Open in YouTube} button is always present, and a link with no
eleven-character video id --- a channel or a playlist --- opens externally rather than pretending to
play.

\subsection{Dictionary}

A glossary entry is a term, an English definition and a Vietnamese meaning, in a dedicated
\code{meaningVi} field rather than one free-text blob, because a reader of English-language
material wants both and wants them separable. An entry belongs to one post or is global.

While reading, the reader shows a chip for every entry attached to this post \emph{or} whose term
appears in the body, matched case-insensitively, so a global term surfaces only on the articles
that use it. Tapping a chip opens a bottom sheet rather than a dialog, so the definition appears
without covering the paragraph the word was read in. An \emph{Open dictionary} link sits beside the
chips for when a definition is not enough.

The \textbf{Dictionary} screen is the glossary as a whole, with full editing: a search box
filtering on term, definition and Vietnamese meaning at once; scope chips for \emph{All},
\emph{Global} and \emph{Per post}; rows showing the term, its definition and its origin (``Global
term'', or ``From \emph{Post title}''). The \code{+} button adds a global term, since anything
created here belongs to no single post. A row tap edits; the row menu edits, jumps to the owning
post, or deletes with an undo snackbar. Adding a term that already exists in the same scope
updates that row instead of stacking a duplicate chip in the reader.

\subsection{Finding things}

\textbf{Search} is one text box, debounced at 300\,ms, plus five filter chips. The box drives an
FTS4 index over post titles and content; because \code{posts\_fts} is content-backed by
\code{nodes}, Room's triggers keep it in step and no application code writes to the index. The
typed text becomes a MATCH query with punctuation stripped and each token given a prefix star, so
``dist sys'' matches \emph{Distributed Systems} while a stray quote cannot break the query.

The result list merges five passes and de-duplicates: FTS hits on title and body, a title
substring pass (so branches and folders match too), posts carrying a matching tag, glossary
entries, and references. Posts sort first, then title-prefix matches, then alphabetically, capped
at fifty; a post matching only in its body carries an excerpt around the first match, snapped to
word boundaries, so the row says why it is there. The chips are \emph{Posts} and \emph{Branches}
(mutually exclusive), \emph{Favorites}, and \emph{Reading}/\emph{Finished} (also mutually
exclusive); set with nothing typed, they list the library on their own. Glossary and reference hits
appear only when no filter is active, having no type, status or star to filter on. Tapping a result
goes where it lives: a post opens the reader, a branch or folder Browse at that node, a glossary hit
its owning post or the dictionary, and a reference the post that cites it.

\textbf{Home} covers the parts you should not have to search for: a greeting chosen by the hour, the
post count, four shortcut tiles (graph, favourites, dictionary, tags), a horizontal \emph{Continue
reading} strip, \emph{Recent posts}, \emph{Favorites}, and a branch list opening Browse at the
branch tapped.

\subsection{The AI assistant}

The \textbf{AI} tab lists chat sessions; a conversation carries four mode chips --- \emph{Answer},
\emph{Suggest}, \emph{Modify} and \emph{Organize} --- deciding the system prompt and what the model
may return. The assistant reads the library only through read-only context queries, so it is never
uploaded wholesale, and it cannot write: proposals arrive as structured actions, and the
\textbf{Review changes} screen shows one row per action to preview, edit, relocate or deselect
before the approved subset is applied. \S\ref{sec:ai} covers the modes, the protocol, the
validation rules and the transaction behaviour in full.

\subsection{Backup, settings and theming}

\textbf{Export} writes the whole library through the system file picker, into a file named for the
current date with a \code{.mocha.json} suffix: one versioned JSON document holding nodes, links,
tags, post--tag pairs, glossary entries and references. Chat history is deliberately not in it ---
a transcript of questions is not part of the knowledge base, and exporting it would put messages
into a file people share.

\textbf{Import} takes a file from the same picker, reads it without writing anything, then asks:
this file holds \emph{n} posts --- \emph{Merge} into the library, or \emph{Replace} everything?
Restore always reassigns primary keys and rewrites every foreign key against the new ids, so a
merge cannot collide with rows already present and a replace cannot inherit stale ones. Nodes go
in parents-first; a node whose parent is missing lands at the top level rather than disappearing,
and a per-post glossary entry whose post is missing is dropped. The search index needs no
attention, because it is content-backed.

Because the device is the only copy, the app nags. A \emph{Remind me weekly} toggle posts a
notification on the next cold start once a week has passed since the last export --- or since first
launch, tracked separately so someone who has never exported is not told they have. The
notification permission is requested when the toggle is switched on, and the check runs at startup
rather than on a schedule, avoiding a background scheduler for a message that only has to be seen
when the app is next opened.

\textbf{Settings} also holds the theme (\emph{Follow system}, \emph{Light}, \emph{Dark}), applied
through \code{AppCompatDelegate} and kept in \code{SharedPreferences} so it can be read
synchronously in \code{Application.onCreate}, before any view inflates and therefore without a
flash of the wrong palette. Below it sits the \textbf{AI gateway address}, with \emph{Save},
\emph{Reset} and \emph{Test connection}; the address is normalised to end in a slash, defaults to
the emulator loopback, and the test reports plainly --- ``Connected to the gateway'', or ``Could not
reach the gateway. Your library still works.'' A privacy paragraph beneath states what leaves the
device and what never does.

\subsection{Offline and degraded states}

Every screen except the assistant works with the network switched off, and works identically ---
none of them makes a network call at all. That follows from the local-first commitment in
\S\ref{sec:intro}; it is not a fallback path that has to be maintained separately.

When the \textbf{gateway is unreachable}, both chat screens ping its health endpoint on entry and
show a banner with a retry button, the send button disabled, worded reassuringly: ``AI unavailable
--- your library still works.'' Nothing is written and nothing is queued, and a reply mid-stream
when the connection dies leaves no trace either (\S\ref{sec:ai}). A YouTube sheet with no network
swaps the player for ``Video unavailable offline. Open it in YouTube to try again.''

\textbf{Empty and error states} are one mechanism rather than thirteen.
\code{ui/common/ListState} is a five-value enum --- \code{LOADING}, \code{CONTENT}, \code{EMPTY},
\code{ERROR}, \code{OFFLINE} --- and \code{ListStateBinder} renders it into a shared overlay of
spinner, message and retry button. All thirteen list-bearing screens bind it, each supplying its
own wording, so an empty state says something useful instead of ``no data'': the graph says
``Nothing to draw yet. Write a few posts and connect them with \code{[[wiki-links]]}''.

One honest gap: the \code{OFFLINE} value is defined in the enum and handled in the binder, but no
ViewModel emits it. The chat screens express unreachability through their own banner and a disabled
send button instead, which is the better fit --- an offline chat should still show its history
rather than be replaced by a full-screen message --- so that branch of the binder is unused code
today rather than a state the user can reach.
